% Part: first-order-logic
% Chapter: beyond
% Section: higher-order-logic

\documentclass[../../../include/open-logic-section]{subfiles}

\begin{document}

\olfileid{fol}{byd}{hol}

\olsection{Logica van een steeds hogere orde}

Met tweede-ordelogica konden we binnen de formele taal spreken over
verzamelingen van objecten uit het eerste-orde!!{domain}. We hoeven daar niet
op te houden. In derde-ordelogica kunnen we bijvoorbeeld werken met
verzamelingen van zulke verzamelingen. We kunnen dan misschien ook
verzamelingen toelaten waarin zowel objecten als verzamelingen objecten zitten.
Met vierde-ordelogica kunnen we weer over verzamelingen van dat soort dingen
spreken. Deze stap kunnen we zo vaak herhalen als we willen.

In de praktijk kiezen logici bij de !!{formula}ring van hogere-ordelogica vaak
voor functies in plaats van relaties. Als we functies en relaties op de
gebruikelijke manier met elkaar vereenzelvigen, maakt dat geen wezenlijk
verschil. Begin met enkele basale ``soorten'' $A$, $B$, $C$,~\dots. Vanaf nu
noemen we zo'n soort een ``type''. Nieuwe typen maken we met de volgende regel:
\begin{quote}
Als $\sigma$ en $\tau$ eindige typen zijn, dan is ook $\sigma \to \tau$ een
eindig type.
\end{quote}
Een type is hier een syntactisch label. Het label zegt wat voor object iets is
en welke plaats het in ons !!{domain} krijgt. Het type $\sigma \to \tau$ hoort
bij functies die een object van type~$\sigma$ als invoer nemen en een object
van type~$\tau$ als uitvoer geven. Neem bijvoorbeeld een type $\Omega$ voor de
waarheidswaarden ``waar'' en ``onwaar'', en een type $\Nat$ voor de natuurlijke
getallen. Een object van type $\Nat \to \Omega$ geeft bij elk natuurlijk getal
een waarheidswaarde. We kunnen zo'n object dus zien als een unaire relatie, of
als een deelverzameling van $\Nat$. Een object van type $\Nat \to \Nat$ is een
functie van natuurlijke getallen naar natuurlijke getallen. Een object van
type $(\Nat \to \Nat) \to \Nat$ neemt zelf zo'n functie als invoer en levert
een natuurlijk getal als uitvoer. Zo'n functie van een hoger type heet een
``functionaal''.

We kunnen hogere-ordelogica zien als meersoortige logica: voor elk type is er
dan een eigen soort. Meestal is het duidelijker om de syntaxis van de logica
met hogere typen rechtstreeks op te bouwen. We definiëren eerst inductief de
eindige typen:
\begin{enumerate}
\item $\Nat$ is een eindig type.
\item Als $\sigma$ en $\tau$ eindige typen zijn, dan is ook $\sigma
  \to \tau$ een eindig type.
\item Als $\sigma$ en $\tau$ eindige typen zijn, dan is ook $\sigma \times
  \tau$ een eindig type.
\end{enumerate}
De bedoeling van deze drie regels is als volgt. $\Nat$ is het type van de
natuurlijke getallen. $\sigma \to \tau$ is het type van functies van
$\sigma$ naar $\tau$. $\sigma \times \tau$ is het type van paren: het eerste
object komt uit $\sigma$ en het tweede uit $\tau$. Daarna definiëren we
inductief welke termen er zijn:
\begin{enumerate}
\item Voor elk type $\sigma$ is er een voorraad variabelen $x$, $y$,
  $z$, \dots van type $\sigma$
\item $\Obj 0$ is een term van type $\Nat$
\item $\Obj S$ (opvolger) is een term van type $\Nat \to \Nat$
\item Als $s$ een term van type $\sigma$ is en $t$ een term van type
  $\Nat \to (\sigma \to \sigma)$, dan is $\Obj{R}_{st}$ een term van type
  $\Nat \to \sigma$
\item Als $s$ een term van type $\tau \to \sigma$ is en $t$ een
  term van type~$\tau$, dan is $s(t)$ een term van type $\sigma$
\item Als $s$ een term van type~$\sigma$ is en $x$ een variabele van
  type~$\tau$, dan is $\lambd[x][s]$ een term van type $\tau \to \sigma$.
\item Als $s$ een term van type~$\sigma$ is en $t$ een term van
  type~$\tau$, dan is $\tuple{s, t}$ een term van type $\sigma \times
  \tau$.
\item Als $s$ een term van type~$\sigma \times \tau$ is, dan is $p_1(s)$ een
  term van type~$\sigma$ en $p_2(s)$ een term van type~$\tau$.
\end{enumerate}
De term $\Obj{R}_{st}$ staat voor een recursief gedefinieerde functie. De
eerste regel geeft de beginwaarde en de tweede regel maakt elke volgende
waarde uit de vorige:
\begin{align*}
\Obj{R}_{st}(0) & = s \\
\Obj{R}_{st}(x+1) & = t(x, R_{st}(x)),
\end{align*}
De term $\tuple{s, t}$ staat voor het paar met $s$ als eerste component en
$t$ als tweede component. De termen $p_1(s)$ en~$p_2(s)$ zijn de
``projecties'': ze halen respectievelijk het eerste en het tweede element uit
$s$. Ten slotte staat $\lambd[x][s]$ voor de functie~$f$ die voldoet aan
\[
f(x) = s
\]
voor elke~$x$ van type~$\sigma$. Regel (6) geeft ons dus een vorm van
comprehensie: we kunnen functies met termen definiëren. !!^{formula}s beginnen
met !!{identity}suitspraken $\eq[s][t]$ tussen termen van hetzelfde type.
Daaruit bouwen we verdere !!^{formula}s met de gewone propositionele
connectieven en met kwantoren voor hogere typen. Als axioma's nemen we de
basisvergelijkingen die het gedrag van de hierboven ingevoerde termen
vastleggen. Daar voegen we de gewone logicaregels voor kwantoren en
!!{identity} aan toe.

We kunnen aan het systeem van eindige typen ook een type $\Omega$ voor
waarheidswaarden toevoegen. Dan hebben we extra axioma's nodig die vastleggen
hoe dat type werkt. Met een handige opzet kunnen we complexe !!{formula}s
zelfs helemaal vervangen door termen van type $\Omega$!{} Daarna passen we het
bewijssysteem daarop aan. Wat zo ontstaat, is in wezen de
\emph{eenvoudige typentheorie} die Alonzo Church in de jaren dertig opstelde.

Ook voor logica met hogere typen kunnen we verschillende semantieken kiezen.
In de volledige versie loopt een variabele van type $\sigma \to \tau$ over
de verzameling van
\emph{alle} functies van objecten van type~$\sigma$ naar objecten van
type~$\tau$. Deze semantiek is zo sterk dat er geen volledig en effectief
!!{derivation} systeem voor bestaat. Er is ook een zwakkere semantiek. Een
!!a{structure} bevat daarin verzamelingen elementen $T_\tau$, één voor elk
type $\tau$. De !!{structure} bevat bovendien de passende bewerkingen voor het
toepassen van functies, het nemen van projecties enzovoort. Als al deze
onderdelen goed worden uitgewerkt, krijgen we volledigheidsstellingen voor de
hierboven beschreven soorten !!{derivation} systemen.

Logica met hogere typen geeft een natuurlijk raamwerk waarin we een groot deel
van de wiskunde kunnen onderbrengen. Als we met $\Nat$ beginnen, kunnen we de
reële getallen en continue functies definiëren. Het raamwerk
past ook bijzonder goed bij intuïtionistische logica. De typen hebben daar een
duidelijke ``constructieve'' uitleg: ze geven aan wat voor constructies als
objecten gelden. We kunnen zelfs constructieve semantieken voor hogere typen
ontwikkelen. Die zijn gebaseerd op intuïtionistische logica in plaats van op
klassieke logica. Ze maken de constructieve uitleg helder en zijn in veel
opzichten interessanter dan de klassieke versies.

\end{document}
